\documentclass[10pt]{beamer}

% --- Theme ---
\usetheme{Madrid}
\usecolortheme{default}
\setbeamertemplate{navigation symbols}{}
\setbeamertemplate{footline}{\hfill\insertframenumber/\inserttotalframenumber\hspace*{5pt}\vspace{2pt}}

% --- Packages ---
\usepackage{amsmath}
\usepackage{booktabs}

% ============================================================
\title{Adaptive EEG-to-Spike Framework}
\subtitle{Brain Wave Classification using Spiking Neural Networks}
\author{
  Bala Mosay J \quad (2023BTIA029) \\[2pt]
  Allwyn Jeffo Raj A \quad (2023BTIA002)
}
\institute{Under the guidance of Dr.\ Vaibhav Meshram}
\date{September 2, 2026}

\begin{document}

% ============================================================
\begin{frame}
\titlepage
\end{frame}

% ============================================================
\begin{frame}{What We Are Building}

\small
\begin{itemize}
  \item A system that reads brain waves and guesses what the person is imagining
  \item The person imagines moving their left hand, right hand, feet, or just resting
  \item We use 64 electrodes on the scalp to capture brain signals
  \item The signals are processed by a \textbf{Spiking Neural Network} (SNN) --- an AI that works like a real brain
  \item SNNs use \textbf{spikes} (tiny pulses) instead of constant numbers, so they use much less energy
  \item This makes them ideal for battery-powered wearable devices
\end{itemize}

\end{frame}

% ============================================================
\begin{frame}{The 6-Stage Pipeline}

\small
\begin{enumerate}
  \item \textbf{Clean the Data} --- Remove noise, eye blinks, and muscle artifacts from raw brain recordings
  \item \textbf{Extract Features} --- Find the most useful patterns in the cleaned signals using CSP (Common Spatial Patterns)
  \item \textbf{Pick Best Frequencies} --- Select the frequency bands (mu: 8--12\,Hz, beta: 13--30\,Hz) that carry the most information
  \item \textbf{Convert to Spikes} --- Turn regular numbers into binary spikes (0 or 1) using a threshold
  \item \textbf{SNN Classifies} --- Feed spikes into a Spiking Neural Network that votes over 16 timesteps
  \item \textbf{Train and Evaluate} --- Teach the network using labeled examples, then test on unseen data
\end{enumerate}

\vspace{4pt}
Each stage is a separate Python module. The full pipeline runs with one command.

\end{frame}

% ============================================================
\begin{frame}{Code: Data Processing}

\small
\textbf{5 files handle all data loading and cleaning:}

\vspace{4pt}
\begin{tabular}{ll}
\toprule
\textbf{File} & \textbf{What It Does} \\
\midrule
\texttt{loader.py} & Downloads brain data from PhysioNet for all 103 subjects \\
\texttt{preprocessing.py} & Applies bandpass filter, ICA, epoching, AutoReject \\
\texttt{events.py} & Labels each trial with the correct class (left/right/feet/rest) \\
\texttt{balancer.py} & Ensures equal number of trials per class \\
\texttt{normalization.py} & Scales data so all subjects are comparable \\
\bottomrule
\end{tabular}

\vspace{4pt}
All 103 subjects' cleaned data is cached in \texttt{outputs/preprocessed/} so preprocessing only runs once.

\end{frame}

% ============================================================
\begin{frame}{Code: Feature Extraction}

\small
\textbf{2 files extract useful patterns from brain signals:}

\vspace{4pt}
\begin{tabular}{ll}
\toprule
\textbf{File} & \textbf{What It Does} \\
\midrule
\texttt{rcsp.py} & Implements Regularized CSP (spatial filter) \\
\texttt{bandpass.py} & Applies CSP per frequency band, computes log-variance \\
\bottomrule
\end{tabular}

\vspace{6pt}
\textbf{The key equation (CSP features):}
\vspace{-2pt}
\[
f_b = \log\!\left(\mathrm{Var}(\text{filtered\_signal}_b) + \epsilon\right)
\]

\vspace{-4pt}
Translation: ``For each frequency band, measure how much the signal varies, then take the logarithm.''

\vspace{4pt}
Result: \textbf{32 features} per trial (2 bands $\times$ 4 classes $\times$ 4 spatial filters).

\end{frame}

% ============================================================
\begin{frame}{Code: SNN Model}

\small
\textbf{2 files build the brain-like AI:}

\vspace{4pt}
\begin{tabular}{ll}
\toprule
\textbf{File} & \textbf{What It Does} \\
\midrule
\texttt{encoder.py} & Converts feature numbers into binary spikes (0/1) \\
\texttt{lif_net.py} & Implements the LIF (Leaky Integrate-and-Fire) network \\
\bottomrule
\end{tabular}

\vspace{6pt}
\textbf{The key equation (LIF neuron):}
\vspace{-2pt}
\[
V[t] = V[t\!-\!1] \times \left(1 - \frac{1}{\tau}\right) + \text{input}[t]
\]

\vspace{-4pt}
Translation: ``A neuron accumulates input in a bucket. The bucket leaks slowly. When it is full, the neuron fires a spike and empties.''

\vspace{4pt}
Architecture: Linear(32$\to$64) $\to$ LIF $\to$ Dropout(0.5) $\to$ Linear(64$\to$32) $\to$ LIF $\to$ Linear(32$\to$4).

\end{frame}

% ============================================================
\begin{frame}{Code: Training \& Evaluation}

\small
\textbf{3 files handle training, testing, and energy measurement:}

\vspace{4pt}
\begin{tabular}{ll}
\toprule
\textbf{File} & \textbf{What It Does} \\
\midrule
\texttt{trainer.py} & Trains the SNN using Adam optimizer + cosine LR schedule \\
\texttt{evaluator.py} & Computes accuracy, kappa, F1 score \\
\texttt{synops.py} & Counts Synaptic Operations (energy proxy) \\
\bottomrule
\end{tabular}

\vspace{6pt}
\textbf{The key equation (SynOps --- energy metric):}
\vspace{-2pt}
\[
\text{SynOps} = \sum_{l=1}^{L} \text{fanout}_l \times \bar{s}_l
\]

\vspace{-4pt}
Translation: ``Count how many neuron connections fire on average. Fewer spikes = less energy.''

\vspace{4pt}
Baseline SynOps: \textbf{0.72 per trial}. Target after Phase 3: lower.

\end{frame}

% ============================================================
\begin{frame}{39 Tests Pass}

\small
\textbf{Every module has unit tests.} We ran all 39 tests --- all pass.

\vspace{4pt}
\begin{tabular}{ll}
\toprule
\textbf{Test File} & \textbf{What It Checks} \\
\midrule
\texttt{test\_preprocessing.py} & Filtering works, correct number of epochs \\
\texttt{test\_rcsp.py} & CSP filters are different across folds (no leakage) \\
\texttt{test\_encoder.py} & Spikes are binary, density near 25\% \\
\texttt{test\_lif\_net.py} & Output shape correct, membrane stays positive \\
\texttt{test\_trainer.py} & Loss decreases during training \\
\texttt{test\_synops.py} & SynOps are non-negative and consistent \\
\texttt{test\_no\_leakage.py} & Train/test parameters never mix \\
\bottomrule
\end{tabular}

\vspace{4pt}
\textbf{No data leakage} is enforced by design: z-score, CSP filters, feature scaling, and threshold are all fit on the training fold only.

\end{frame}

% ============================================================
\begin{frame}{Results So Far}

\small
\textbf{Subject-mixed trial-level evaluation (103 subjects, 5-fold CV):}

\vspace{6pt}
\begin{tabular}{lr}
\toprule
\textbf{Classifier} & \textbf{Accuracy} \\
\midrule
Chance (random guess) & 25.0\% \\
Our SNN (baseline) & \textbf{31.8\%} $\pm$ 1.1 \\
LogReg (same features) & 34.6\% $\pm$ 0.8 \\
\bottomrule
\end{tabular}

\vspace{6pt}
\textbf{Key finding:} The SNN (31.8\%) is only 2.8\% behind LogReg (34.6\%) on the same features. The model is not the bottleneck --- the features are.

\vspace{4pt}
Within a single subject, LogReg achieves \textbf{42.0\%} --- the features do work. This
subject-mixed comparison is not a leave-one-subject-out result.

\end{frame}

% ============================================================
\begin{frame}{Bottleneck Analysis}

\small
\textbf{Why is accuracy low?}

\vspace{4pt}
\begin{itemize}
  \item Subject-mixed accuracy (34.6\%) is much lower than within-subject (42.0\%)
  \item Different people's brains produce different patterns for the same task
  \item Fixed frequency bands (mu: 8--12\,Hz, beta: 13--30\,Hz) don't work equally well for everyone
  \item The CSP features appear subject-specific; LOSO is needed to confirm transfer
\end{itemize}

\vspace{4pt}
\textbf{This is exactly what Phase 3 addresses:}

\vspace{4pt}
\begin{itemize}
  \item \textbf{C1 (Adaptive Bands):} Pick the best frequencies for each person
  \item \textbf{C2 (Learnable Threshold):} Let the AI find the optimal spike threshold
  \item \textbf{C3 (PLIF Network):} Use fast and slow neurons to capture different time scales
\end{itemize}

\vspace{4pt}
Target: 55--65\% accuracy with lower energy.

\end{frame}

% ============================================================
\begin{frame}{Thank You}

\begin{center}
\vspace{10pt}
{\Large\textbf{Thank You}}\\[12pt]
{\large Questions \& Discussion}

\vspace{20pt}
\textbf{Team:}\\
Bala Mosay J \quad (2023BTIA029)\\
Allwyn Jeffo Raj A \quad (2023BTIA002)

\vspace{10pt}
\textbf{Guide:}\\
Dr.\ Vaibhav Meshram
\end{center}

\end{frame}

\end{document}
